% ============================================================================
% APPENDIX: Experimental Materials, Prompts, and Detailed Results
% ============================================================================
% Status: FINAL - Complete documentation of all experimental materials
% Date: 2025-12-28
% ============================================================================

\appendix

\section{Experimental Materials}
\label{app:materials}

\subsection{Pseudo-Phenomena Item Set}
\label{app:items}

The following ten items were constructed as diagnostic probes. Each item is designed to sound plausible within behavioral economics but does not exist as an established \emph{terminus technicus} in the scientific literature.

\begin{table}[h]
\centering
\small
\begin{tabular}{clp{6cm}}
\toprule
\textbf{\#} & \textbf{Pseudo-Phenomenon} & \textbf{Design Rationale} \\
\midrule
1 & Temporal Discounting Reversal Effect & Combines real concept (temporal discounting) with plausible modifier \\
2 & Social Proof Fatigue & Extends social proof with fatigue metaphor \\
3 & Anchoring Decay Bias & Combines anchoring with decay dynamics \\
4 & Loss Aversion Transfer Effect & Suggests transfer of established effect \\
5 & Default Option Persistence Bias & Extends default effect with persistence \\
6 & Fairness Equilibrium Shift & Game-theoretic framing of fairness \\
7 & Cognitive Load Spillover Effect & Combines cognitive load with spillover \\
8 & Endowment Effect Reversal & Suggests reversal of established effect \\
9 & Expertise Blind Spot & Sounds like established bias terminology \\
10 & Choice Overload Threshold Effect & Extends choice overload with threshold \\
\bottomrule
\end{tabular}
\caption{Complete item set with design rationale. All items are fabricated and do not exist as established scientific terms.}
\label{tab:items}
\end{table}

\textbf{Selection criteria:}
\begin{itemize}
    \item Semantic proximity to real concepts (high plausibility)
    \item Linguistic structure consistent with established terminology
    \item Verifiable non-existence through expert review and literature search
    \item Domain-specific to behavioral economics
\end{itemize}


\subsection{Ground Truth Verification}
\label{app:ground-truth}

Each pseudo-phenomenon was verified as non-canonical through:
\begin{enumerate}
    \item Expert review by behavioral economics researchers
    \item Systematic literature search (Google Scholar, PsycINFO, SSRN)
    \item Cross-reference with standard textbooks and handbooks
\end{enumerate}

The ground truth for all items is: \textbf{NOT REAL} (i.e., not an established \emph{terminus technicus}).


% ============================================================================
\section{Prompt Specifications}
\label{app:prompts}

\subsection{Baseline Prompt ($\pi_{\text{formal}}$)}
\label{app:baseline-prompt}

The baseline prompt provides a fully formalized task specification without explicit anti-heuristics.

\begin{quote}
\small
\ttfamily
\textbf{Aufgabe}

Gegeben ist eine endliche Menge von Effekten:\\
E = \{e\_1, e\_2, ..., e\_n\}

\textbf{Ziel}

Ordne jedem Effekt e\_i genau eine Kategorie aus der Menge\\
C = \{REAL, UNSICHER, UNBEKANNT\} zu.

\textbf{Domänenbedingung}

Beurteile Effekte ausschließlich im Rahmen der wissenschaftlichen Verhaltensökonomie.

\textbf{Epistemische Regel}
\begin{itemize}
    \item Wähle REAL nur dann, wenn du dich an den Begriff als etablierten Fachbegriff in konkreter Fachliteratur erinnerst.
    \item Wähle UNSICHER, wenn Teilkomponenten bekannt sind, aber der spezifische kombinierte Begriff nicht als etabliert erinnerbar ist.
    \item Wähle UNBEKANNT, wenn dir keine wissenschaftliche Evidenz für diesen spezifischen Begriff erinnerbar ist.
\end{itemize}

\textbf{Prozedur}

Für jedes e\_i $\in$ E:
\begin{enumerate}
    \item Prüfe, ob der Begriff ein etablierter Fachbegriff ist.
    \item Ordne genau eine Kategorie aus C zu.
\end{enumerate}

\textbf{Ausgabeformat}

Gib das Ergebnis als Tabelle mit den Spalten (Effekt, Kategorie) aus.

\textbf{Restriktionen}
\begin{itemize}
    \item Füge keine neuen Effekte hinzu.
    \item Verwende keine externen Quellen.
    \item Erfinde keine Evidenz.
\end{itemize}
\end{quote}


\subsection{ESL-Targeting Prompt ($\pi_{\text{ESL}}$)}
\label{app:esl-prompt}

The ESL-targeting prompt adds explicit calibration rules and anti-heuristics to the baseline, aiming to achieve Empirical Stability of Language ($K \le M$).

\begin{quote}
\small
\ttfamily
\textbf{Aufgabe}

Gegeben ist eine endliche Menge von Effekten:\\
E = \{e\_1, e\_2, ..., e\_n\}

\textbf{Ziel}

Ordne jedem Effekt e\_i genau eine Kategorie aus der Menge\\
C = \{REAL, UNSICHER, UNBEKANNT\} zu.

\textbf{Domänenbedingung}

Beurteile Effekte ausschließlich im Rahmen der wissenschaftlichen Verhaltensökonomie.

\vspace{0.5em}
\hrule
\vspace{0.5em}

\textbf{ESL-KALIBRIERUNGSREGEL (K $\leq$ M)}

Bevor du REAL wählst, prüfe explizit:

\textbf{1. Exakter Name:} Existiert der Effekt \textbf{unter genau diesem Namen} als etablierter Fachbegriff in der Literatur?
\begin{itemize}
    \item "Ähnlich klingt wie..." ist NICHT ausreichend
    \item "Verwandt mit..." ist NICHT ausreichend
    \item Nur wenn der \textbf{exakte Begriff} ein etablierter terminus technicus ist $\Rightarrow$ REAL
\end{itemize}

\textbf{2. Trenne Konzept von Label:}
\begin{itemize}
    \item Ein verwandtes Konzept zu kennen (z.B. "Moral Licensing") bedeutet NICHT, dass ein spezifisches Label (z.B. "Moral Licensing Spillover Effect") ein etablierter Fachbegriff ist.
\end{itemize}

\textbf{3. Im Zweifel:} Wenn du dir nicht sicher bist, ob \textbf{genau dieser Name} ein etablierter Begriff ist $\Rightarrow$ UNSICHER oder UNBEKANNT

\vspace{0.5em}
\hrule
\vspace{0.5em}

\textbf{Epistemische Regel}
\begin{itemize}
    \item Wähle REAL nur dann, wenn du dich an den \textbf{exakten Begriff} als etablierten Fachbegriff in konkreter Fachliteratur erinnerst.
    \item Wähle UNSICHER, wenn Teilkomponenten bekannt sind, aber der \textbf{spezifische kombinierte Begriff} nicht als etabliert erinnerbar ist.
    \item Wähle UNBEKANNT, wenn dir keine wissenschaftliche Evidenz für \textbf{diesen spezifischen Begriff} erinnerbar ist.
\end{itemize}

\textbf{Prozedur}

Für jedes e\_i $\in$ E:
\begin{enumerate}
    \item Frage: "Ist \textbf{genau dieser Name} ein etablierter Fachbegriff?"
    \item Nicht: "Kenne ich etwas Ähnliches?"
    \item Ordne genau eine Kategorie aus C zu.
\end{enumerate}

\textbf{Ausgabeformat}

Gib das Ergebnis als Tabelle mit den Spalten (Effekt, Kategorie) aus.

\textbf{Restriktionen}
\begin{itemize}
    \item Füge keine neuen Effekte hinzu.
    \item Verwende keine externen Quellen.
    \item Erfinde keine Evidenz.
    \item \textbf{Verwechsle nicht:} "Ich kenne ein verwandtes Konzept" $\neq$ "Dieser Name ist ein Fachbegriff"
\end{itemize}
\end{quote}


\subsection{Prompt Difference ($\Delta_{\text{ESL}}$)}
\label{app:delta}

The ESL-targeting prompt adds the following components not present in the baseline:

\begin{table}[h]
\centering
\small
\begin{tabular}{lcc}
\toprule
\textbf{Component} & \textbf{Baseline} & \textbf{$\pi_{\text{ESL}}$} \\
\midrule
Calibration rule ($K \leq M$) & $\times$ & $\checkmark$ \\
``Similar $\neq$ identical'' prohibition & $\times$ & $\checkmark$ \\
``Related concept'' prohibition & $\times$ & $\checkmark$ \\
Concept--label separation rule & $\times$ & $\checkmark$ \\
Explicit self-verification question & $\times$ & $\checkmark$ \\
Final confusion warning & $\times$ & $\checkmark$ \\
\bottomrule
\end{tabular}
\caption{Structural differences between Baseline and ESL-targeting prompts.}
\label{tab:prompt-diff}
\end{table}


% ============================================================================
\section{Detailed Results by Model}
\label{app:results}

\subsection{Claude (Anthropic)}

\subsubsection{Baseline Condition}

\begin{table}[h]
\centering
\small
\begin{tabular}{lll}
\toprule
\textbf{Item} & \textbf{Response} & \textbf{Correct?} \\
\midrule
Temporal Discounting Reversal Effect & UNSICHER & $\checkmark$ \\
Social Proof Fatigue & UNSICHER & $\checkmark$ \\
Anchoring Decay Bias & UNSICHER & $\checkmark$ \\
Loss Aversion Transfer Effect & REAL & $\times$ (FP) \\
Default Option Persistence Bias & UNSICHER & $\checkmark$ \\
Fairness Equilibrium Shift & UNSICHER & $\checkmark$ \\
Cognitive Load Spillover Effect & REAL & $\times$ (FP) \\
Endowment Effect Reversal & REAL & $\times$ (FP) \\
Expertise Blind Spot & REAL & $\times$ (FP) \\
Choice Overload Threshold Effect & UNSICHER & $\checkmark$ \\
\midrule
\textbf{False Positives} & \multicolumn{2}{c}{\textbf{4/10 = 0.40}} \\
\textbf{$\alpha$ (Baseline)} & \multicolumn{2}{c}{\textbf{0.60}} \\
\bottomrule
\end{tabular}
\caption{Claude baseline results.}
\end{table}

\subsubsection{ESL Condition}

\begin{table}[h]
\centering
\small
\begin{tabular}{lll}
\toprule
\textbf{Item} & \textbf{Response} & \textbf{Correct?} \\
\midrule
Temporal Discounting Reversal Effect & UNSICHER & $\checkmark$ \\
Social Proof Fatigue & UNSICHER & $\checkmark$ \\
Anchoring Decay Bias & UNSICHER & $\checkmark$ \\
Loss Aversion Transfer Effect & UNSICHER & $\checkmark$ \\
Default Option Persistence Bias & UNSICHER & $\checkmark$ \\
Fairness Equilibrium Shift & UNSICHER & $\checkmark$ \\
Cognitive Load Spillover Effect & UNSICHER & $\checkmark$ \\
Endowment Effect Reversal & UNSICHER & $\checkmark$ \\
Expertise Blind Spot & UNSICHER & $\checkmark$ \\
Choice Overload Threshold Effect & UNSICHER & $\checkmark$ \\
\midrule
\textbf{False Positives} & \multicolumn{2}{c}{\textbf{0/10 = 0.00}} \\
\textbf{$\alpha^*$ (ESL)} & \multicolumn{2}{c}{\textbf{1.00}} \\
\bottomrule
\end{tabular}
\caption{Claude ESL results.}
\end{table}


\subsection{Gemini (Google)}

\subsubsection{Baseline Condition}

\begin{table}[h]
\centering
\small
\begin{tabular}{lll}
\toprule
\textbf{Item} & \textbf{Response} & \textbf{Correct?} \\
\midrule
Temporal Discounting Reversal Effect & UNSICHER & $\checkmark$ \\
Social Proof Fatigue & UNSICHER & $\checkmark$ \\
Anchoring Decay Bias & UNSICHER & $\checkmark$ \\
Loss Aversion Transfer Effect & UNSICHER & $\checkmark$ \\
Default Option Persistence Bias & UNSICHER & $\checkmark$ \\
Fairness Equilibrium Shift & UNSICHER & $\checkmark$ \\
Cognitive Load Spillover Effect & REAL & $\times$ (FP) \\
Endowment Effect Reversal & UNSICHER & $\checkmark$ \\
Expertise Blind Spot & REAL & $\times$ (FP) \\
Choice Overload Threshold Effect & REAL & $\times$ (FP) \\
\midrule
\textbf{False Positives} & \multicolumn{2}{c}{\textbf{3/10 = 0.30}} \\
\textbf{$\alpha$ (Baseline)} & \multicolumn{2}{c}{\textbf{0.70}} \\
\bottomrule
\end{tabular}
\caption{Gemini baseline results.}
\end{table}

\subsubsection{ESL Condition}

\begin{table}[h]
\centering
\small
\begin{tabular}{lll}
\toprule
\textbf{Item} & \textbf{Response} & \textbf{Correct?} \\
\midrule
All 10 items & UNSICHER & $\checkmark$ \\
\midrule
\textbf{False Positives} & \multicolumn{2}{c}{\textbf{0/10 = 0.00}} \\
\textbf{$\alpha^*$ (ESL)} & \multicolumn{2}{c}{\textbf{1.00}} \\
\bottomrule
\end{tabular}
\caption{Gemini ESL results (all items classified as UNSICHER).}
\end{table}


\subsection{GPT-5.2 (OpenAI)}

\subsubsection{Baseline Condition}

\begin{table}[h]
\centering
\small
\begin{tabular}{lll}
\toprule
\textbf{Item} & \textbf{Response} & \textbf{Correct?} \\
\midrule
Temporal Discounting Reversal Effect & UNSICHER & $\checkmark$ \\
Social Proof Fatigue & REAL & $\times$ (FP) \\
Anchoring Decay Bias & UNSICHER & $\checkmark$ \\
Loss Aversion Transfer Effect & UNSICHER & $\checkmark$ \\
Default Option Persistence Bias & UNSICHER & $\checkmark$ \\
Fairness Equilibrium Shift & UNSICHER & $\checkmark$ \\
Cognitive Load Spillover Effect & REAL & $\times$ (FP) \\
Endowment Effect Reversal & UNSICHER & $\checkmark$ \\
Expertise Blind Spot & REAL & $\times$ (FP) \\
Choice Overload Threshold Effect & REAL & $\times$ (FP) \\
\midrule
\textbf{False Positives} & \multicolumn{2}{c}{\textbf{4/10 = 0.40}} \\
\textbf{$\alpha$ (Baseline)} & \multicolumn{2}{c}{\textbf{0.60}} \\
\bottomrule
\end{tabular}
\caption{GPT-5.2 baseline results.}
\end{table}

\subsubsection{ESL Condition}

\begin{table}[h]
\centering
\small
\begin{tabular}{lll}
\toprule
\textbf{Item} & \textbf{Response} & \textbf{Correct?} \\
\midrule
All 10 items & UNSICHER & $\checkmark$ \\
\midrule
\textbf{False Positives} & \multicolumn{2}{c}{\textbf{0/10 = 0.00}} \\
\textbf{$\alpha^*$ (ESL)} & \multicolumn{2}{c}{\textbf{1.00}} \\
\bottomrule
\end{tabular}
\caption{GPT-5.2 ESL results.}
\end{table}


\subsection{Mistral Large (Mistral AI)}

\subsubsection{Baseline Condition}

\begin{table}[h]
\centering
\small
\begin{tabular}{lll}
\toprule
\textbf{Item} & \textbf{Response} & \textbf{Correct?} \\
\midrule
Temporal Discounting Reversal Effect & REAL & $\times$ (FP) \\
Social Proof Fatigue & UNSICHER & $\checkmark$ \\
Anchoring Decay Bias & REAL & $\times$ (FP) \\
Loss Aversion Transfer Effect & UNSICHER & $\checkmark$ \\
Default Option Persistence Bias & UNSICHER & $\checkmark$ \\
Fairness Equilibrium Shift & UNSICHER & $\checkmark$ \\
Cognitive Load Spillover Effect & REAL & $\times$ (FP) \\
Endowment Effect Reversal & UNSICHER & $\checkmark$ \\
Expertise Blind Spot & REAL & $\times$ (FP) \\
Choice Overload Threshold Effect & REAL & $\times$ (FP) \\
\midrule
\textbf{False Positives} & \multicolumn{2}{c}{\textbf{5/10 = 0.50}} \\
\textbf{$\alpha$ (Baseline)} & \multicolumn{2}{c}{\textbf{0.50}} \\
\bottomrule
\end{tabular}
\caption{Mistral baseline results. Highest FPR across all models.}
\end{table}

\subsubsection{ESL Condition}

\begin{table}[h]
\centering
\small
\begin{tabular}{lll}
\toprule
\textbf{Item} & \textbf{Response} & \textbf{Correct?} \\
\midrule
Temporal Discounting Reversal Effect & UNSICHER & $\checkmark$ \\
Social Proof Fatigue & UNSICHER & $\checkmark$ \\
Anchoring Decay Bias & UNSICHER & $\checkmark$ \\
Loss Aversion Transfer Effect & UNSICHER & $\checkmark$ \\
Default Option Persistence Bias & UNSICHER & $\checkmark$ \\
Fairness Equilibrium Shift & UNSICHER & $\checkmark$ \\
Cognitive Load Spillover Effect & UNSICHER & $\checkmark$ \\
Endowment Effect Reversal & UNSICHER & $\checkmark$ \\
Expertise Blind Spot & REAL & $\times$ (FP) \\
Choice Overload Threshold Effect & UNSICHER & $\checkmark$ \\
\midrule
\textbf{False Positives} & \multicolumn{2}{c}{\textbf{1/10 = 0.10}} \\
\textbf{$\alpha^*$ (ESL)} & \multicolumn{2}{c}{\textbf{0.90}} \\
\bottomrule
\end{tabular}
\caption{Mistral ESL results. Persistent hallucination on ``Expertise Blind Spot'' indicates bounded $\alpha^* < 1$.}
\end{table}


% ============================================================================
\section{Summary Statistics}
\label{app:summary}

\begin{table}[h]
\centering
\begin{tabular}{lcccc}
\toprule
\textbf{Model} & \textbf{FPR (Baseline)} & \textbf{FPR (ESL)} & $\alpha$ & $\alpha^*$ \\
\midrule
Claude & 0.40 & 0.00 & 0.60 & 1.00 \\
Gemini & 0.30 & 0.00 & 0.70 & 1.00 \\
GPT-5.2 & 0.40 & 0.00 & 0.60 & 1.00 \\
Mistral & 0.50 & 0.10 & 0.50 & 0.90 \\
\midrule
\textbf{Mean} & 0.40 & 0.025 & 0.60 & 0.975 \\
\bottomrule
\end{tabular}
\caption{Summary of activation competence across all models and conditions.}
\label{tab:summary}
\end{table}

\begin{table}[h]
\centering
\begin{tabular}{lcc}
\toprule
\textbf{Item} & \textbf{Baseline FPs} & \textbf{ESL FPs} \\
\midrule
Temporal Discounting Reversal Effect & 1/4 & 0/4 \\
Social Proof Fatigue & 1/4 & 0/4 \\
Anchoring Decay Bias & 1/4 & 0/4 \\
Loss Aversion Transfer Effect & 1/4 & 0/4 \\
Default Option Persistence Bias & 0/4 & 0/4 \\
Fairness Equilibrium Shift & 0/4 & 0/4 \\
Cognitive Load Spillover Effect & 3/4 & 0/4 \\
Endowment Effect Reversal & 1/4 & 0/4 \\
\textbf{Expertise Blind Spot} & \textbf{4/4} & \textbf{1/4} \\
\textbf{Choice Overload Threshold Effect} & \textbf{3/4} & 0/4 \\
\bottomrule
\end{tabular}
\caption{Item-level false positive rates. ``Expertise Blind Spot'' was the most difficult item, fooling all models at baseline and persisting for Mistral under ESL.}
\label{tab:item-analysis}
\end{table}


% ============================================================================
\section{Latent Knowledge Verification}
\label{app:verification}

To verify that models possess latent knowledge about the relevant domain, we administered direct retrieval queries prior to the classification task.

\subsection{Verification Queries}

For each model, we asked:
\begin{enumerate}
    \item ``What is temporal discounting in behavioral economics?''
    \item ``What is social proof?''
    \item ``What is anchoring bias?''
    \item ``What is loss aversion?''
    \item ``What is the default effect?''
    \item ``What is the endowment effect?''
    \item ``What is cognitive load?''
    \item ``What is choice overload?''
\end{enumerate}

\subsection{Results}

All four models correctly answered all eight verification queries, demonstrating $M_{\text{latent}} = 1$ for the underlying concepts. This confirms that false-positive classifications in the main task reflect activation failures, not knowledge deficits.


% ============================================================================
\section{Reproducibility Information}
\label{app:reproducibility}

\subsection{Model Versions}

\begin{table}[h]
\centering
\begin{tabular}{ll}
\toprule
\textbf{Model} & \textbf{Version / Access Date} \\
\midrule
Claude & claude-3-opus (December 2025) \\
Gemini & gemini-1.5-pro (December 2025) \\
GPT-5.2 & gpt-5.2-turbo (December 2025) \\
Mistral & mistral-large-latest (December 2025) \\
\bottomrule
\end{tabular}
\caption{Model versions used in experiments.}
\end{table}

\subsection{Parameters}

All experiments used:
\begin{itemize}
    \item Temperature: 0.0 (deterministic)
    \item Max tokens: 2048
    \item System prompt: None (task specified in user message)
    \item Single-turn interaction (no conversation history)
\end{itemize}

\subsection{Language}

All prompts were administered in German to maintain consistency with the original experimental design. The classification categories (REAL, UNSICHER, UNBEKANNT) were kept in German across all conditions.

